\documentclass[12pt]{article}
\usepackage{amsmath}
\usepackage{amssymb}
\usepackage{geometry}
\usepackage{listings}
\usepackage{xcolor}
\usepackage{hyperref}
\usepackage{enumitem}
\usepackage{booktabs}

\geometry{margin=1in}

\title{Architectural Selection Debate: MLP, LSTM, and CNN1D for DDoS Detection}
\subtitle{Analyzing Model Choices for Tabular Network Flow Data Classification}
\author{}
\date{}

\begin{document}

\maketitle

\section{Problem Context}

The task is \textbf{binary classification} of network flow data for DDoS attack detection:
\begin{itemize}
    \item \textbf{Data Type:} Tabular/structured data (network flow features)
    \item \textbf{Features:} Packet count, byte count, duration, flows, packet rate, etc.
    \item \textbf{Task:} Classify each sample as Normal (0) or DDoS Attack (1)
    \item \textbf{Challenge:} Choosing appropriate neural network architecture for tabular data
\end{itemize}

\section{The Fundamental Question}

\textbf{Why use MLP, LSTM, and CNN1D when the data is tabular, not sequential or image-based?}

This document presents arguments for and against each architecture choice.

\section{MLP (Multi-Layer Perceptron)}

\subsection{Arguments FOR MLP}

\subsubsection{1. Natural Fit for Tabular Data}
\begin{itemize}
    \item \textbf{Designed for structured data:} MLPs are the standard architecture for tabular/structured data
    \item \textbf{No assumptions:} Doesn't assume spatial or temporal structure in features
    \item \textbf{Feature interactions:} Fully connected layers can model any feature interaction
    \item \textbf{Interpretability:} Easier to understand and debug than CNNs/LSTMs
\end{itemize}

\subsubsection{2. Theoretical Justification}
\[
\hat{y} = \text{softmax}(W_2 \cdot \text{ReLU}(W_1 \cdot x + b_1) + b_2)
\]
\begin{itemize}
    \item Universal function approximator (can approximate any continuous function)
    \item No architectural constraints on feature relationships
    \item Direct mapping from features to output
\end{itemize}

\subsubsection{3. Computational Efficiency}
\begin{itemize}
    \item \textbf{Fast training:} Parallel processing, no sequential dependencies
    \item \textbf{Fewer parameters:} More parameter-efficient than CNN/LSTM
    \item \textbf{Lower memory:} Simpler architecture requires less memory
    \item \textbf{Scalability:} Scales well with data size
\end{itemize}

\subsubsection{4. Modern Enhancements}
\begin{itemize}
    \item \textbf{BatchNormalization:} Stabilizes training, allows higher learning rates
    \item \textbf{Dropout:} Prevents overfitting effectively
    \item \textbf{Regularization:} Well-regularized MLPs are competitive with complex models
\end{itemize}

\subsection{Arguments AGAINST MLP}

\subsubsection{1. Limited Hierarchical Learning}
\begin{itemize}
    \item \textbf{No feature abstraction:} May not learn hierarchical feature representations
    \item \textbf{Feature engineering dependent:} Relies heavily on good feature engineering
    \item \textbf{Flat structure:} All features treated equally, no multi-scale patterns
\end{itemize}

\subsubsection{2. Curse of Dimensionality}
\begin{itemize}
    \item Can struggle with very high-dimensional feature spaces
    \item May not handle sparse features well
    \item Requires careful feature selection
\end{itemize}

\subsection{MLP Verdict}

\textbf{✅ BEST CHOICE} for tabular data classification. Theoretically sound, computationally efficient, and appropriate for the problem domain.

\section{LSTM (Long Short-Term Memory)}

\subsection{Arguments FOR LSTM}

\subsubsection{1. Sequential Feature Processing}
\begin{itemize}
    \item \textbf{Sequential processing:} Processes features sequentially, potentially learning dependencies
    \item \textbf{Memory mechanism:} Can "remember" important features while processing others
    \item \textbf{Stateful learning:} Maintains hidden state that could encode feature relationships
\end{itemize}

\subsubsection{2. Temporal Interpretation}
\begin{itemize}
    \item \textbf{Network flow as temporal:} Network flows have temporal characteristics (packets arrive over time)
    \item \textbf{Feature ordering:} If features are ordered by importance or temporal relevance, LSTM can leverage this
    \item \textbf{Dependency modeling:} Can model complex dependencies between features through hidden states
\end{itemize}

\subsubsection{3. Empirical Success}
\begin{itemize}
    \item LSTMs have been used successfully in network intrusion detection
    \item Can model complex decision boundaries
    \item Proven in security domain applications
\end{itemize}

\subsubsection{4. Mathematical Formulation}
\[
\begin{aligned}
h_t &= \text{LSTM}(x_t, h_{t-1}, c_{t-1}) \\
\hat{y} &= \text{softmax}(W \cdot h_T + b)
\end{aligned}
\]
\begin{itemize}
    \item Processes features as sequence: $x_1, x_2, \ldots, x_T$
    \item Maintains memory through hidden state $h_t$ and cell state $c_t$
    \item Final hidden state $h_T$ captures sequential information
\end{itemize}

\subsection{Arguments AGAINST LSTM}

\subsubsection{1. No True Temporal Structure}
\begin{itemize}
    \item \textbf{Tabular data is IID:} Each sample is independent—no temporal dependencies between samples
    \item \textbf{Feature vector $\neq$ time series:} Treating features as timesteps is conceptually incorrect
    \item \textbf{Misuse of architecture:} LSTM is designed for sequences where order matters (time series, text), not feature vectors
\end{itemize}

\subsubsection{2. Computational Overhead}
\begin{itemize}
    \item \textbf{More parameters:} LSTM has significantly more parameters than MLP
    \begin{align*}
    \text{LSTM params} &= 4 \times (d \times d + d \times d_{in} + d) \\
    &\text{where } d = \text{units}, d_{in} = \text{input size}
    \end{align*}
    \item \textbf{Slower training:} Sequential processing is slower than parallel MLP processing
    \item \textbf{Vanishing gradients:} Can suffer from gradient issues in deep networks
\end{itemize}

\subsubsection{3. No Clear Benefit}
\begin{itemize}
    \item \textbf{Feature relationships:} MLPs can learn feature interactions through fully connected layers
    \item \textbf{Attention mechanisms:} Modern tabular models (e.g., TabNet) use attention, not RNNs
    \item \textbf{Over-engineering:} More complex than necessary for tabular data
\end{itemize}

\subsubsection{4. Feature Order Dependency}
\begin{itemize}
    \item \textbf{Order sensitivity:} Changing feature order changes LSTM behavior
    \item \textbf{Arbitrary ordering:} Features like `port_no`, `protocol`, `tx_kbps` don't have inherent ordering
    \item \textbf{Inconsistent results:} Different feature orders may yield different results
\end{itemize}

\subsection{LSTM Verdict}

\textbf{⚠️ QUESTIONABLE CHOICE} for tabular data. Over-engineered for the problem, but may work empirically if features have meaningful ordering.

\section{CNN1D (1D Convolutional Neural Network)}

\subsection{Arguments FOR CNN1D}

\subsubsection{1. Feature Pattern Recognition}
\begin{itemize}
    \item \textbf{Local patterns:} CNN1D can detect local patterns and correlations between adjacent features
    \item \textbf{Example:} Features like `pktcount`, `bytecount`, `duration_sec` might have meaningful relationships when positioned together
    \item \textbf{Translation invariance:} CNN learns patterns regardless of exact feature positions (though less relevant for tabular data)
\end{itemize}

\subsubsection{2. Hierarchical Feature Learning}
\begin{itemize}
    \item \textbf{Multi-scale patterns:} Multiple Conv1D layers with pooling can learn hierarchical representations
    \item \textbf{Feature abstraction:} Lower layers detect simple patterns (e.g., "high packet count"), higher layers detect complex combinations
    \item \textbf{Reduced overfitting:} Convolutional layers share weights, reducing parameters compared to fully connected layers
\end{itemize}

\subsubsection{3. Mathematical Formulation}
\[
y[i] = \sum_{j=0}^{k-1} w[j] \cdot x[i+j] + b
\]
\begin{itemize}
    \item \textbf{Convolution operation:} Detects patterns in local windows
    \item \textbf{Multiple filters:} Each filter learns different patterns
    \item \textbf{Feature maps:} Output represents detected patterns at each position
\end{itemize}

\subsubsection{4. Empirical Success}
\begin{itemize}
    \item CNNs have shown success in tabular data tasks when features are treated as sequences
    \item Can capture non-linear interactions between features more efficiently than MLPs
    \item Proven in various classification tasks
\end{itemize}

\subsubsection{5. Architecture Benefits}
\begin{itemize}
    \item \textbf{MaxPooling:} Extracts dominant features, reduces dimensionality
    \item \textbf{Parameter sharing:} Same filter used across all positions
    \item \textbf{Multi-layer abstraction:} Simple patterns → complex patterns → classification
\end{itemize}

\subsection{Arguments AGAINST CNN1D}

\subsubsection{1. No Spatial Structure}
\begin{itemize}
    \item \textbf{Tabular data lacks spatial/temporal ordering:} Features like `port_no`, `protocol`, `tx_kbps` don't have inherent spatial relationships
    \item \textbf{Arbitrary feature order:} Changing feature order shouldn't affect results, but CNNs are sensitive to input order
    \item \textbf{Misaligned assumptions:} CNNs assume local patterns matter, but in tabular data, any feature can interact with any other
\end{itemize}

\subsubsection{2. Over-engineering}
\begin{itemize}
    \item \textbf{Unnecessary complexity:} MLPs are designed for tabular data and may be more appropriate
    \item \textbf{Parameter efficiency:} CNN1D might not be more efficient than MLP for this problem size
    \item \textbf{Feature engineering:} Still requires good feature engineering
\end{itemize}

\subsubsection{3. Data Transformation Artifact}
\begin{itemize}
    \item The `expand_dims(axis=2)` operation artificially creates a "channel" dimension
    \item This is a workaround, not a natural representation
    \item Features are not truly sequential or spatial
\end{itemize}

\subsubsection{4. Limited Receptive Field}
\begin{itemize}
    \item \textbf{Small kernel size:} Kernel size 3 only sees 3 adjacent features
    \item \textbf{Long-range dependencies:} May miss relationships between distant features
    \item \textbf{Requires depth:} Needs multiple layers to capture global patterns
\end{itemize}

\subsection{CNN1D Verdict}

\textbf{⭐ EXPERIMENTAL CHOICE} for tabular data. May capture local feature patterns, but lacks theoretical justification for non-spatial data.

\section{Comparative Analysis}

\subsection{Architectural Comparison}

\begin{table}[h]
\centering
\begin{tabular}{|l|l|l|l|}
\hline
\textbf{Aspect} & \textbf{MLP} & \textbf{LSTM} & \textbf{CNN1D} \\
\hline
\textbf{Data Assumption} & None (IID) & Sequential & Spatial/Local \\
\hline
\textbf{Feature Order} & Independent & Dependent & Dependent \\
\hline
\textbf{Parameters} & Low & High & Medium \\
\hline
\textbf{Training Speed} & Fast & Slow & Medium \\
\hline
\textbf{Interpretability} & High & Medium & Low \\
\hline
\textbf{Theoretical Fit} & ✅ Excellent & ❌ Poor & ⚠️ Questionable \\
\hline
\textbf{Empirical Fit} & ✅ Good & ⚠️ Variable & ⚠️ Variable \\
\hline
\end{tabular}
\caption{Architectural Comparison}
\end{table}

\subsection{Parameter Count Comparison}

For $d_{\text{features}} = 20$ and binary classification:

\begin{align*}
\text{MLP:} &\quad \sim 3,500 \text{ parameters} \\
\text{LSTM:} &\quad \sim 29,400 \text{ parameters} \\
\text{CNN1D:} &\quad \sim 49,600 \text{ parameters}
\end{align*}

\textbf{Conclusion:} MLP is most parameter-efficient.

\subsection{Computational Complexity}

\begin{align*}
\text{MLP:} &\quad O(d_{\text{features}} \times d_{\text{hidden}}) \quad \text{(parallel)} \\
\text{LSTM:} &\quad O(T \times d^2) \quad \text{(sequential, } T = d_{\text{features}}) \\
\text{CNN1D:} &\quad O(T \times k \times F) \quad \text{(convolution, } k = \text{kernel}, F = \text{filters})
\end{align*}

\textbf{Conclusion:} MLP has lowest computational complexity.

\section{Why All Three Models Were Chosen}

\subsection{Research Methodology}

The selection of MLP, LSTM, and CNN1D represents an \textbf{experimental approach}:

\begin{enumerate}
    \item \textbf{Baseline (MLP):} Provides theoretically sound baseline for comparison
    \item \textbf{Experimental (LSTM/CNN1D):} Tests whether sequential/spatial assumptions help despite being tabular data
    \item \textbf{Comparative Analysis:} Enables empirical comparison of different architectural paradigms
\end{enumerate}

\subsection{Valid Research Questions}

\begin{itemize}
    \item \textbf{Q1:} Can sequential processing (LSTM) improve tabular classification?
    \item \textbf{Q2:} Can local pattern detection (CNN1D) capture feature relationships?
    \item \textbf{Q3:} Do complex architectures outperform simple MLP for this problem?
    \item \textbf{Q4:} Is the added complexity justified by performance gains?
\end{itemize}

\subsection{Expected Outcomes}

\begin{itemize}
    \item \textbf{MLP:} Expected to perform well (baseline)
    \item \textbf{LSTM:} May work if features have meaningful ordering
    \item \textbf{CNN1D:} May work if local feature patterns exist
    \item \textbf{Best Model:} Determined empirically through experiments
\end{itemize}

\section{Recommendations}

\subsection{For This Problem}

\begin{enumerate}
    \item \textbf{Start with MLP} as primary model (theoretically sound)
    \item \textbf{Use LSTM/CNN1D} for comparison and ablation studies
    \item \textbf{Focus on feature engineering} rather than complex architectures
    \item \textbf{Consider alternatives:} TabNet, XGBoost, or Gradient Boosting might be more appropriate
\end{enumerate}

\subsection{For Future Work}

\begin{itemize}
    \item \textbf{Feature importance analysis:} Determine if feature ordering matters
    \item \textbf{Architecture ablation:} Systematically test which components help
    \item \textbf{Hybrid approaches:} Combine MLP with attention mechanisms
    \item \textbf{Ensemble methods:} Combine predictions from all three models
\end{itemize}

\section{Conclusion}

\subsection{Summary}

\begin{itemize}
    \item \textbf{MLP:} ✅ Theoretically sound, appropriate baseline
    \item \textbf{LSTM:} ⚠️ Questionable for tabular data, but worth testing
    \item \textbf{CNN1D:} ⚠️ Experimental, may capture local patterns
\end{itemize}

\subsection{Final Verdict}

The architectural choices reflect \textbf{valid research methodology}:
\begin{itemize}
    \item Testing whether architectures designed for other domains can be adapted to tabular data
    \item Empirical results should determine which model performs best
    \item Theoretical justification is weak for LSTM/CNN1D, but empirical validation is key
\end{itemize}

\textbf{Key Insight:} While MLP is theoretically optimal for tabular data, LSTM and CNN1D serve as experimental baselines to test if sequential/spatial assumptions provide empirical benefits, even without strong theoretical foundation.

\section{Mathematical Formulations}

\subsection{MLP Forward Pass}

\[
\begin{aligned}
h_1 &= \text{ReLU}(W_1 \cdot x + b_1) \\
h_2 &= \text{ReLU}(W_2 \cdot h_1 + b_2) \\
\hat{y} &= \text{softmax}(W_3 \cdot h_2 + b_3)
\end{aligned}
\]

\subsection{LSTM Forward Pass}

\[
\begin{aligned}
f_t &= \sigma(W_f \cdot [h_{t-1}, x_t] + b_f) \\
i_t &= \sigma(W_i \cdot [h_{t-1}, x_t] + b_i) \\
o_t &= \sigma(W_o \cdot [h_{t-1}, x_t] + b_o) \\
\tilde{c}_t &= \tanh(W_C \cdot [h_{t-1}, x_t] + b_C) \\
c_t &= f_t \odot c_{t-1} + i_t \odot \tilde{c}_t \\
h_t &= o_t \odot \tanh(c_t) \\
\hat{y} &= \text{softmax}(W \cdot h_T + b)
\end{aligned}
\]

\subsection{CNN1D Forward Pass}

\[
\begin{aligned}
y_1[i] &= \sum_{j=0}^{k-1} w_1[j] \cdot x[i+j] + b_1 \quad \text{(Conv1D Layer 1)} \\
p_1[i] &= \max(y_1[2i], y_1[2i+1]) \quad \text{(MaxPooling)} \\
y_2[i] &= \sum_{j=0}^{k-1} w_2[j] \cdot p_1[i+j] + b_2 \quad \text{(Conv1D Layer 2)} \\
p_2[i] &= \max(y_2[2i], y_2[2i+1]) \quad \text{(MaxPooling)} \\
z &= \text{Flatten}(p_2) \\
h &= \text{ReLU}(W_1 \cdot z + b_1) \\
\hat{y} &= \text{softmax}(W_2 \cdot h + b_2)
\end{aligned}
\]

\section{References and Further Reading}

\begin{itemize}
    \item MLP: Universal function approximators for tabular data
    \item LSTM: Originally designed for sequential data (time series, text)
    \item CNN1D: Designed for 1D signals (audio, time series, sequences)
    \item Tabular Data: Traditional ML (XGBoost, Random Forest) often outperforms deep learning
\end{itemize}

\end{document}

